\title{LongRoPE: Extending LLM Context Window Beyond 2 Million Tokens}

\begin{abstract}
The ability to process large context windows is a sought-after characteristic in large language models (LLMs). Nevertheless, extending context windows beyond approximately 128k tokens remains challenging, primarily because of significant fine-tuning expenses, limited availability of lengthy textual data, and the detrimental impact of extreme values associated with novel token positions.
\end{abstract}

This work presents LongRoPE, marking the first demonstration of extending the context window of pre-trained LLMs up to a remarkable \textbf{2048k} tokens, accomplished with as few as 1k fine-tuning steps at training lengths no greater than 256k, all while preserving the original short context performance. The proposed method hinges on three primary innovations: (i) it recognizes and leverages two distinct non-uniformities in positional interpolation through an efficient search procedure, leading to superior initialization for fine-tuning and facilitating an 8$\times$ extension without fine-tuning; (ii) it implements a staged extension approach by initially fine-tuning an LLM at 256k length and subsequently applying a second positional interpolation to the already extended LLM, thus realizing a 2048k context window; (iii) it reoptimizes LongRoPE for the 8k sequence length to recoup short context performance. Comprehensive experiments conducted on LLaMA2 and Mistral over diverse benchmarks confirm the efficacy of the approach. The LongRoPE methodology preserves the original model architecture except for minimal changes to the positional embedding, allowing compatibility with nearly all existing optimizations. The code will be released at \texttt{https://github.com/microsoft/LongRoPE}

\section{Introduction}

Large Language Models (LLMs) have achieved impressive results across a range of tasks~\cite{gpt4,llama2}, yet they are constrained by a finite context window; for example, LLaMA2 is restricted to 4096 tokens~\cite{llama2}. When input exceeds this context length, the effectiveness of LLMs deteriorates, as the models have not been trained on such extended input positions. This limitation creates difficulties in key applications, such as in-context learning with a substantial number of exemplars~\cite{Huang2023FewerIM} and in operating LLM agents~\cite{park2023generative,madaan2023selfrefine}.

Recent studies demonstrate that fine-tuning a pre-trained LLM enables the context window to be pushed to approximately 128k tokens~\cite{longlora,pi,yarn,zhang2024soaring,liu2023scaling}. Nonetheless, several significant challenges limit further expansions of the context window. Firstly, incorporating previously unseen positional indices during extension introduces a substantial number of problematic or aberrant values, which can cause out-of-distribution complications and hamper the stability and convergence of fine-tuning~\cite{pi}. This problem becomes especially pronounced when scaling from 4k up to over 1000k, as this process can involve over 90\% new position indices. Secondly, effective fine-tuning generally depends on the availability of texts that match the desired longer lengths. Unfortunately, present datasets offer limited instances of extremely lengthy documents, particularly those surpassing 1000k tokens. In addition, training on such ultra-long sequences demands considerable computational resources, significantly escalating both the duration of training and GPU usage to impractical levels. Thirdly, when scaling the context to such great lengths, the model's attention mechanism tends to become overly distributed across many tokens, which diminishes its capacity to perform well on tasks involving the original, shorter contexts~\cite{pi}.

A common strategy to address the initial challenge involves interpolating the RoPE positional embeddings~\cite{rope,pi} by compressing the indices of new positions to align with the pretrained range, as illustrated in Fig.\ref{fig:overview}. Position Interpolation (PI)~\cite{pi} employs a linear interpolation of RoPE’s rotary angles based on the proportion of extension. NTK~\cite{ntk,dynamicntk}, on the other hand, proposes non-uniform methods for interpolation and extrapolation across different dimensions of RoPE. YaRN~\cite{yarn} further subdivides RoPE dimensions into three frequency-centric categories, applying extrapolation, NTK, or linear interpolation to each group respectively. Despite these innovations, positional embeddings in the Transformer structure possess a \emph{complex and heterogeneous distribution of information entropy}. Current methodologies fail to capitalize on this nuanced irregularity, resulting in information degradation and constraining the maximum achievable context window.

Section~\ref{sec:analysis} identifies two principal empirical insights: \textbf{(1)} Successful positional interpolation necessitates accounting for two types of non-uniformity—variation across RoPE dimensions and across token positions. Reduced interpolation is advantageous for lower RoPE dimensions and for tokens located near the start, although the ideal configuration is contingent upon the specific extended length being targeted. \textbf{(2)} Incorporating these non-uniformities into the design of positional interpolation preserves more of the information embedded in the original RoPE, especially within key dimensions and crucial token positions. This approach reduces the informational loss inherent in positional interpolation, resulting in superior initial conditions for subsequent fine-tuning. Furthermore, it supports extending context by a factor of 8$\times$ even without the need for fine-tuning.

Building on these results, we propose \textbf{LongRoPE}, a robust approach that extends the context window of LLMs to more than 2 \emph{million} tokens. LongRoPE's design centers on three primary innovations. First, {LongRoPE} leverages multidimensional, non-uniform characteristics in positional interpolation to their fullest extent. It determines optimal rescale factors for the rotation angles in each dimension of RoPE by considering token positions. Given that the search space for suitable rescale factors expands exponentially with increasing extension ratios, {LongRoPE} applies an evolutionary search strategy coupled with two optimization methods to enhance the efficiency of the search process. An illustration of the rescaled RoPE derived from this search is provided in Fig.~\ref{fig:overview}.

Subsequently, {LongRoPE} adopts a resource-efficient, stepwise extension method to obtain a 2048k context window, bypassing the need for direct fine-tuning on ultra-long text sequences, which are both rare and difficult to source. This process starts by identifying a viable 256k sequence length for the pre-trained LLM, followed by fine-tuning at this scale. Next, leveraging our non-uniform positional interpolation—which permits an eightfold ($8\times$) increase in context length without additional fine-tuning—we perform a further search for optimal RoPE rescaling parameters using the already extended and fine-tuned LLM. Through this two-stage approach, the 2048k context window is successfully realized for both LLaMA2 and Mistral~\cite{mistral}.

In order to address the potential drop in performance on the original, shorter context lengths, {LongRoPE} maintains the practice of modifying the RoPE rescale factors for the extended LLM. In an analogous fashion to scaling up the context window from 256k to 2048k, we apply our search algorithm to reduce the context length to 4k and 8k on the LLM fine-tuned at 256k, aiming to reduce the degree of positional interpolation. At inference time, if the input sequence is shorter than 8k tokens, we substitute the original RoPE with the identified rescale factors.

A comprehensive set of experiments utilizing multiple LLMs and a range of long-context tasks validates the efficacy of our proposed approach. Our findings indicate that LongRoPE consistently sustains low perplexity across evaluation lengths from 4k up to 2048k, attains more than 90\% accuracy in passkey retrieval tasks, and achieves accuracy comparable to leading models on established benchmarks created within a 4096-token context window. The LongRoPE technique is compatible with all LLMs utilizing RoPE embeddings. We intend to make our codebase and the LongRoPE-2048k models publicly available.

\section{Non-uniformity in Positional Interpolation}
\label{sec:analysis}

\subsection{Preliminary}

Transformer models require explicit positional information, often in the form of position embedding, to represent the order of input tokens. Our work focuses on the RoPE~\cite{rope} position embedding, which is widely used in recent  LLMs. For a token at position index $n$, its corresponding RoPE encoding can be simplified as follows:

\begin{equation}
		\fontsize{7.8}{7.8} \selectfont
	\label{eq:rope}
	\left[ cos(n\theta_0), sin(n\theta_0), cos(n\theta_1), \cdot\cdot\cdot, cos(n\theta_{d/2-1}), sin(n\theta_{d/2-1}) \right]   
\end{equation}

where $d$ denotes the dimension of the embedding space, 
$n\theta_i$ refers to the rotary angle assigned to the token at position $n$, 
and $\theta_i=\theta^{-2i/d}$ encodes the corresponding rotation frequencies. RoPE typically utilizes 10000 as the default base value for $\theta$.

\textbf{\emph{Context window extension ratio $s$} and positional interpolation}. We introduce $s$ as the quotient of the expanded context length $L'$ and the original context length $L$: $s=\frac{L'}{L}$.

To extend the context window from $L$ to $L'$, current positional interpolation methods suggest downscaling rotation frequencies $\theta_i$ based on extension ratio $s$. Let $\beta=\theta^{2/d}$, and $\lambda$ denote the actual rescale factor related to $s$, we   unify these positional interpolation methods as follows:

\begin{equation}	
	\fontsize{7.5}{7.5} \selectfont
	\label{eq:unified_rope}
	\begin{aligned}
		\left[cos\left(\frac{n}{\lambda(\beta)^0}\right), sin \left(\frac{n}{\lambda(\beta)^0}\right), cos\left(\frac{n}{\lambda(\beta)^1}\right), \cdot\cdot\cdot,  sin \left(\frac{n}{\lambda(\beta)^{d/2-1}}\right) \right]   
	\end{aligned}
\end{equation}

\textbf{Linear positional interpolation (PI)}. PI~\cite{pi} employs a linear interpolation strategy for position indices that operates within the original pre-training sequence length. When extending to a new target length by a ratio $s$, the rotation angles assigned to each position are uniformly scaled down by a factor of $\lambda=s$ for all RoPE dimensions. This scaling effectively compresses position information, leading to increased overlap and making it more difficult for the model to differentiate between tokens that are near each other. Consequently, at higher extension ratios, PI typically exhibits reduced performance.

\textbf{NTK-based interpolation and extrapolation}. ~\cite{ntk,dynamicntk} examine RoPE through the lens of information encoding, leveraging the Neural Tangent Kernel (NTK) theory~\cite{jacot2018neural,tancik2020fourier}. To address the problem of dense positional information in PI, they advocate for spreading interpolation demands over the various dimensions of RoPE. This approach involves applying a smaller scaling factor to the lower (higher frequency) dimensions while assigning greater scaling to the higher (lower frequency) dimensions, thereby enabling both interpolation and extrapolation for positional encodings, with $\lambda=s^{i}$. The advanced dynamic NTK technique~\cite{dynamicntk} further customizes the extension ratio for each position by considering the sequence length at hand. Distinct from PI, which requires additional fine-tuning, NTK-informed approaches allow for context window enlargement without fine-tuning; however, they are typically limited to a maximum scaling ratio of 4$\times$.

\textbf{YaRN}~\cite{yarn} presents a notable advancement in positional interpolation by segmenting the RoPE dimensions into three groups according to frequency and assigning distinct interpolation methods to each group. Specifically, the high-frequency dimensions are subjected to extrapolation ($\lambda$=1), whereas the low-frequency ones leverage linear interpolation (PI). Dimensions situated between these two extremes are processed using the NTK approach. The central innovation of YaRN is its frequency-based partitioning of RoPE dimensions; however, this assignment presently relies on manual empirical tuning, which can lead to less-than-optimal outcomes when applied to newly developed LLMs.

\subsection{Study on Non-uniform Positional Interpolation}

Drawing inspiration from NTK and YaRN, we observe that their improvements stem from introducing non-linearity—most notably, through differential treatment of frequencies along RoPE dimensions to facilitate more targeted interpolation and extrapolation. Nonetheless, prevailing non-linear mechanisms depend significantly on heuristics and human-crafted approaches. This observation prompts two central questions: (1) Does the present method for positional interpolation represent the best possible choice? (2) Might there exist additional, as-yet undiscovered non-linear strategies?

In addressing these questions, we employ evolution search (refer to Sec.~\ref{sec:method}) to identify improved non-uniform positional interpolation strategies for LLaMA2-7B. This process is steered by perplexity measurements, utilizing 5 randomly selected instances from the PG19~\cite{pg19} validation dataset. Our experiments yield the following principal observations.

\textbf{Finding 1}: \textit{RoPE dimensions exhibit substantial non-uniformities, which are not effectively handled by current positional interpolation methods.}

We optimize $\lambda$ individually for each RoPE dimension as defined in Eq.~\ref{eq:unified_rope}. Table~\ref{tbl:exp1} presents a comparison of LLaMA2-7B perplexity across various approaches on the PG19 and Proof-pile~\cite{proof-pile} test sets, all evaluated in the absence of fine-tuning. The optimized parameters obtained through our search yield notable perplexity reductions, indicating that prevailing interpolation strategies—both linear (PI) and non-uniform (Dynamic-NTK, YaRN)—fail to achieve optimality. Particularly, YaRN demonstrates inferior performance relative to PI and NTK on PG19, likely because it does not fully cover the intended context window when the LLM is not fine-tuned. Illustratively, with an 8k context, YaRN’s perplexity sharply increases beyond 7k tokens.

Our investigation reveals that the rescaled factors $\lambda$ in Eq.~\ref{eq:unified_rope} exhibit non-uniformity, in contrast to the constant scaling factor $s$ utilized in PI, NTK's formula, and the group-wise computation adopted by YaRN. The introduction of these non-uniform scaling factors yields a marked enhancement in LLaMA2's language modeling ability (as measured by perplexity) for extended context windows of 8k and 16k, even in the absence of fine-tuning. This improvement arises because the resulting positional encodings more faithfully retain the characteristics of the original RoPE, with an emphasis on the preservation of key dimensions, thereby alleviating the challenge faced by LLMs in discriminating among nearby token positions.

\textbf{Finding 2}: \textit{RoPE for the initial tokens in the input sequence should be extrapolated with less interpolation.} 

We propose that for the initial $\hat{n}$ tokens in the input sequence, their RoPE embeddings should involve reduced interpolation. This stems from the observation that these early tokens typically obtain high attention weights, rendering them especially influential in the attention mechanism, as reported by Streaming LLM~\cite{streamingllm} and LM-Infinite~\cite{han2023lminfinite}. To test this proposition, we expand the context window to 8k and 16k using PI and NTK, selectively applying interpolation only beyond the first $\hat n$ tokens, where $\hat n$ is varied (0,2,...,256). Setting $\hat n$=0 recapitulates standard PI and NTK procedures. As demonstrated in Table~\ref{tbl:tokenposition}, two main findings emerge: \textbf{(1)} Excluding the initial tokens from interpolation leads to enhanced model performance; \textbf{(2)} The most effective value of $\hat n$ is contingent on the length of the context window extension.

\textbf{Finding 3}: \textit{Non-uniform positional interpolation effectively extends LLM context window in both fine-tuning and non-fine-tuning settings.} 

Although our results demonstrate that employing the searched non-uniform position interpolation notably enhances out-of-the-box extension capabilities at 8k and 16k context lengths, extending beyond these ranges necessitates additional fine-tuning. Accordingly, we fine-tune LLaMA2-7B using our optimized RoPE method for a context window of 64k tokens (details provided in the Appendix). Table~\ref{tbl:finetune} indicates that our approach yields marked improvements over both PI and YaRN methods, irrespective of whether LLaMA2-7B is pre- or post-fine-tuning. This advantage arises from the effective incorporation of non-uniform positional interpolation, which mitigates information degradation and enables more favorable initialization during fine-tuning.

\textbf{Summary}. This work identifies two key forms of non-uniformity: heterogeneous RoPE dimensions and diverse token positions. Strategic exploitation of these non-uniform characteristics within positional interpolation substantially advances LLMs’ performance for extended context scenarios.

\section{LongRoPE}

\label{sec:method}
Building upon these observations, we propose LongRoPE. This method initially employs an effective search algorithm designed to leverage both non-uniformities to their fullest extent, and subsequently applies this approach to expand the LLM context window to more than 2 million tokens.

\subsection{Problem Formulation}

These two sources of non-uniformity may result in an extensive set of possible solutions and increase the difficulty of the optimization process. To manage this, we reformulate the multidimensional non-uniform position interpolation optimization issue into a search problem.

For a LLM targeting a  context window size of $L'$ and lengthy input documents $\mathbf{X}$, where each $\mathbf{x}\in\mathbf{X}$ surpasses $L'$ in token length, we denote the original rotary angle of the $i^{th}$ dimension in RoPE embedding at token position $n$ as  $\frac{n}{\beta^i}$. The optimization problem is then formulated as follows:

\begin{equation}
\begin{gathered}
		\label{eq:problem}

\underset {\mathbf{x}\in\mathbf{X}; \, |\mathbf{x}|\ge L'}{ \operatorname {arg\,min} } \,  \mathcal{L}\, \left( \text{LLM} (\text{RoPE}, \mathbf{X})\right), \quad \text{where} 
\\
\scalebox{0.93}{
$\underset {\begin{subarray}{l} i=0,\ldots,\frac{d}{2}-1;\\ n\in[0,|\mathbf{x}|);\end{subarray}}{\text{RoPE}_{(n)}}= {\left[\ldots, \cos\left(\mathbb{I}(\hat{\lambda}_{i}, \hat{n}) \frac{n}{\beta^i}\right),\ \sin\left(\mathbb{I}(\hat{\lambda}_{i}, \hat{n}) \frac{n}{\beta^i}\right), \ldots\right] }$}\\
	\text{with} \;\; \mathbb{I}(\hat{\lambda}_{i},\hat{n})=
\begin{cases}
	1 & \text{if}~n<\hat{n} \\
	\frac{1}{\lambda_i} & \text{if}~n \geq \hat{n}
\end{cases}
	\end{gathered}
\end{equation}

In this approach, we define a collection of rescale factors, $\mathbb{I}(\hat{\lambda}_{i},\hat{n})$, to address two types of non-uniformities. Here, $\hat{\lambda}_{i}$ represents the variation across RoPE dimensions, while $\hat{n}$ indicates the threshold for non-uniformity based on token positions. Specifically, $\mathbb{I}(\hat{\lambda}_{i},\hat{n})$ is employed to adjust the rotation angle associated with the $i^{th}$ RoPE dimension. The parameter $\hat{\lambda}_{i}$ functions as the scaling factor, and $\hat{n}$ serves as the positional threshold. For the initial $\hat{n}-1$ token positions, the system retains the original RoPE rotary angle $\frac{n}{\beta^i}$, with the rescaling factor $\hat{\lambda}_{i}$ inactive. Once token positions reach $n \ge \hat{n}$, the rescaling becomes operative.

Suppose the desired context window length is $L'$. The task is to determine the best rescale factors ($\mathbb{I}(\hat{\lambda}_{0},\hat{n})$, $\mathbb{I}(\hat{\lambda}_{1},\hat{n})$, ... $\mathbb{I}(\hat{\lambda}_{i},\hat{n})$, ...) across RoPE dimensions from the $1^{st}$ up to the $d^{th}$. This aims to ensure that the modified $\text{LLM}$, incorporating the updated $\text{RoPE}$, attains the lowest possible next-token prediction loss, $\mathcal{L}$ (reflecting perplexity), when given input samples $\mathbf{X}$ of length $L'$.

\subsection{Searching the Non-uniform Position Interpolation}

In addressing the issue outlined in Eq.~\ref{eq:problem}, we present a straightforward and powerful approach that identifies the most suitable RoPE rescale factors, thereby leveraging the diverse multidimensional variations present in position embedding to their fullest extent.

\textbf{Search space}. To encompass the two types of non-uniformities, we construct an extensive search space. This space is depicted in Table~\ref{tbl:searchspace}. More precisely, we enable the optimization of an individual rescale factor for each dimension within the RoPE embedding. For simplicity, our parameterization involves searching over $\lambda_i$ and $\hat{n}$, rather than directly optimizing $\mathbb{I}(\hat{\lambda}_{i},\hat{n})$, where $\hat{\lambda}_{i}=1/\lambda_i$. As described in Table~\ref{tbl:searchspace}, $\lambda_i$ is permitted to range from a lower bound of 1.0 (which corresponds to direct extrapolation) up to a maximum of $s\times1.25$ (representing interpolation beyond PI), incremented in steps of 0.01. Here, $s$ denotes the ratio by which the context window is extended.

The parameter $\hat{n}$ determines how many of the initial token positions preserve their original positional encoding, foregoing interpolation (specifically, by maintaining the native RoPE embedding). In our experiments, we vary $\hat{n}$ over the set \{0, 1, 2, 4, 8, 12, 16, 20, 24, 28, 32, 64, 128, 256\}. Setting $\hat{n}=0$ results in every token position utilizing the rescale factors identified during the search process.

\textbf{Evolution-based search}. The search space outlined in Table~\ref{tbl:searchspace} encompasses a wide range of options for positional interpolation, making systematic exploration computationally demanding. For instance, a $s=4\times$ scaling results in $400^{128/2}\times14 = 4\times10^{167}$ possible configurations. As the extension ratio increases, the size of the search space grows at an exponential rate. To effectively navigate this complexity, we employ evolution search~\cite{guo2020single} and propose two additional optimization strategies to substantially enhance the efficiency of the search process. The overarching procedure is depicted in Algorithm~\ref{alg:search}.

\textit{Optimized initial population generation}. Rather than fully randomizing the initial population of $P$ rescale factors, we explicitly include the three RoPE rescale factors associated with PI, NTK, and YaRN among the initial individuals. The rest of the $P-3$ individuals in the population are generated by applying random mutations to these three rescale factors with a mutation probability of $p$.

\textit{Monotonically non-decreasing constraint}. Following the generation of the initial population, we evaluate the LLM perplexity for every candidate. In particular, we introduce the relevant RoPE rescale factors into the target LLM and determine the perplexity for the given input $\mathbf{X}$. The individuals ranked in the top $k$ are selected as parents for the evolutionary process. Nevertheless, due to the vastness of the search space, unguided mutation and crossover can frequently produce suboptimal solutions, resulting in redundant perplexity evaluations. This inefficiency becomes more pronounced as $L'$ increases, since each perplexity computation requires extensive inference time.

To resolve this, we enforce a non-decreasing monotonicity condition on the sampled RoPE rescaling parameters, such that $\lambda_i \le \lambda_{i+1}$. Only RoPE rescalings that adhere to this monotonicity criterion are utilized during perplexity evaluation on LLMs, thereby significantly reducing the computational search overhead. In detail, we stipulate that $\lambda_i$ must monotonically increase across the RoPE dimensions ($i=0,...,63$). The imposition of monotonicity over dimensions is grounded in the principles of NTK theory~\cite{jacot2018neural,tancik2020fourier,ntk}, which indicates that dimensions associated with higher frequencies (i.e., lower indices) necessitate less interpolation and hence smaller $\lambda_i$, whereas dimensions with lower frequencies (i.e., higher indices) allow for greater interpolation and thus permit larger $\lambda_i$.

\begin{algorithm}[t]
	\small
	\caption{The search algorithm}
	\textbf{Input:} target LLM, input samples $\mathbf{X}$, population size $P$, mutation size $N_1$, crossover size $N_2$, max iterations $\mathcal{T}$, mutation probability $p$\\

	\begin{algorithmic}[1]
		\label{alg:search}
		\STATE Initialize Top-k as an empty set;	
		\STATE Set $\text{P}_0$ = \textit{Initialize\_population\_with\_optimization}($P$, $\mathbf{X}$, $p$); 
		\FOR{$i = 1$ to $\mathcal{T}$}
		\STATE \textit{Compute\_perplexity} (LLM, $\text{P}_{i-1}$, $\mathbf{X}$);
		\STATE  Update Top-k by calling \textit{Update\_Topk}(Top-k, $\text{P}_{i-1}$);
		\STATE Set $\text{P}_{mutation}$ = \textit{Mutation\_with\_mono\_constraint}(Top-k, $N_1$, $p$); 
		\STATE Generate $\text{P}_{crossover}$ using \textit{Crossover\_with\_mono\_constraint}(Top-k, $N_2$);
		\STATE Form $\text{P}_i$ by taking the union of $\text{P}_{mutation}$, $\text{P}_{crossover}$, and Top-k;
		\ENDFOR
		\STATE Output the member in Top-k with the minimal perplexity;
	\end{algorithmic}
\end{algorithm}

\textbf{8$\times$ extension without fine-tuning}. Utilizing evolutionary search, our approach efficiently discovers non-uniform RoPE rescaling coefficients, selectively maintaining critical dimensions and positional information to reduce information degradation caused by interpolation. As illustrated in Fig.\ref{fig:non-ft}, this strategy successfully increases LLaMA2’s context window from 4k to 32k tokens without the need for fine-tuning. In comparison, previous approaches—including PI, non-uniform NTK, and YaRN—result in a sharp rise in perplexity when the window is expanded beyond 2$\times$.

\subsection{Extending LLM Context Window to 2048K}
\label{sec:extendto2048k}

\textbf{Progressive extension to 2048k}. In this section, we present our approach for enlarging the context window of pre-trained LLMs beyond the standard 4k, reaching up to 2048k. As shown previously, our technique of non-uniform positional interpolation allows for an 8$\times$ increase in context length without the need for fine-tuning. However, achieving more substantial extensions (e.g., 512$\times$) necessitates fine-tuning. One possible strategy involves identifying appropriate RoPE rescaling factors for the intended 2048k context window and subsequently fine-tuning the model. Nonetheless, this path is hindered by the significant computational demands required for training. Furthermore, our practical observations indicate that fine-tuning LLMs for such a substantial extension is highly challenging (refer to Appendix).

Fortunately, LongRoPE performs well not only on the base model but also on LLMs that have undergone further fine-tuning for extended context. To this end, we present a streamlined, incremental approach that can reach a maximum context length of 2048k, requiring only 1k fine-tuning iterations and utilizing a training sequence length capped at 256k.

$\Diamond$ \textit{LongRoPE-based search for scaling pre-trained LLM context to 256k.} Using LLaMA2 as a reference point, we perform a search aimed at achieving context window sizes of 128k and 256k tokens. At this step, the context window is expanded by factors of 32$\times$ and 64$\times$, accordingly.

$\Diamond$ \textit{Fine-tuning to 256k}. In this stage, we adapt the pre-trained LLM to handle a 256k context window. The procedure begins with fine-tuning LLaMA2 for 400 steps, employing RoPE rescaled factors corresponding to 128k. After completing these steps, we update the model by substituting the RoPE rescaled factors for those aligned with 256k, and subsequently carry out 600 more fine-tuning steps. This approach demonstrates greater efficiency compared to direct fine-tuning at 256k from the outset.

$\Diamond$ \textit{Scaling fine-tuned extended LLM to 2048k via LongRoPE search}. In the final step, we conduct an additional search utilizing the LLM that has already been fine-tuned on a 256k context length. This procedure enables us to achieve an exceptionally large context window of 2048k, all without the need for extra fine-tuning. Consequently, this results in a total extension ratio of $512\times$.

\textbf{Recovery in shorter context windows}. Upon expanding the context window to a very large 2048k size, we observe diminished performance inside the initial context window boundaries. This challenge is well-documented for positional interpolation~\cite{pi}, since embedding positions from the high-dimensional original context window into a compressed subset of space causes suboptimal model behavior. Notably, when using a 512$\times$ extension factor, tokens corresponding to the original 4k context window are densely packed, exacerbating this issue.

To address this issue, we conduct an additional evolution search on the extended LLM, specifically tuning the RoPE rescale factors for shorter context lengths such as 4k and 8k. The maximum permissible searched $\lambda$ is decreased, reflecting the reduced necessity for positional interpolation at these lengths. At inference time, the LLM adapts the associated RoPE rescale factors on the fly.

\section{LongRoPE}

\label{sec:method}
Building on these insights, we propose LongRoPE. This approach initially implements an effective search strategy designed to leverage both types of non-uniformity, and subsequently applies this technique to expand the context window of LLMs to exceed 2 million tokens.

\subsection{Problem Formulation}

These two sources of non-uniformity significantly expand the solution space and complicate the optimization process. To tackle this issue, we conceptualize the optimization of multidimensional non-uniform position interpolation as a search problem.

For a LLM targeting a  context window size of $L'$ and lengthy input documents $\mathbf{X}$, where each $\mathbf{x}\in\mathbf{X}$ surpasses $L'$ in token length, we denote the original rotary angle of the $i^{th}$ dimension in RoPE embedding at token position $n$ as  $\frac{n}{\beta^i}$. The optimization problem is then formulated as follows:

\begin{equation}
\begin{gathered}
		\label{eq:problem}

\underset {\mathbf{x}\in\mathbf{X};\, |\mathbf{x}|\ge L'}{ \operatorname {arg\,min} } \, \mathcal{L} \left( \text{LLM} (\text{RoPE}, \mathbf{X}) \right), \;\text{with} \\
\scalebox{0.93}{
$\underset {\begin{subarray}{l} i=0,\ldots,\frac{d}{2}-1;\\ n\in[0,|\mathbf{x}|);\end{subarray}}{ \text{RoPE}_\text{(n)} } = {\left[\cdots, \cos\left( \mathbb{I}(\hat{\lambda}_{i}, \hat{n}) \cdot \frac{n}{\beta^i}\right), \sin\left( \mathbb{I}(\hat{\lambda}_{i}, \hat{n}) \cdot \frac{n}{\beta^i}\right), \cdots \right]}$} \\
\text{where} \; \mathbb{I}(\hat{\lambda}_{i}, \hat{n}) = \begin{cases}
1 & \text{if}~n < \hat{n} \\
{\frac{1}{\lambda}_{i}} & \text{if}~n \geq \hat{n}
\end{cases}
\end{gathered}
\end{equation}

In this context, we define a collection of scaling factors, $\mathbb{I}(\hat{\lambda}_{i},\hat{n})$, designed to address both categories of non-uniformity. Here, $\hat{\lambda}_{i}$ and $\hat{n}$ represent non-uniformities related to the RoPE dimension indices and token positional indices, respectively. More precisely, $\mathbb{I}(\hat{\lambda}_{i},\hat{n})$ is utilized to adjust the rotation angle in the $i^{th}$ RoPE dimension, where $\hat{\lambda}_{i}$ functions as the scaling coefficient and $\hat{n}$ serves as the position threshold for tokens. For the initial $\hat{n}-1$ tokens, the adjustment via $\hat{\lambda}_{i}$ is not employed; thus, the standard RoPE rotary angle $\frac{n}{\beta^i}$ remains unchanged. Only for tokens at positions where $n \ge \hat{n}$ does the rescaling come into effect.

Assuming a desired context window length of $L'$, our task involves identifying the set of optimal rescale factors ($\mathbb{I}(\hat{\lambda}_{0},\hat{n})$, $\mathbb{I}(\hat{\lambda}_{1},\hat{n})$, ...$\mathbb{I}(\hat{\lambda}_{i},\hat{n})$...) across all RoPE dimensions from $1^{st}$ through $d^{th}$. In doing so, the target $\text{LLM}$—now employing these rescaled $\text{RoPE}$ parameters—aims to minimize next-token prediction loss, $\mathcal{L}$ (that is, perplexity), over input sequences $\mathbf{X}$ of length $L'$.

\subsection{Searching the Non-uniform Position Interpolation}

To address the issue presented in Eq.~\ref{eq:problem}, we propose a straightforward but robust approach that identifies optimal RoPE rescale factors, enabling full utilization of the complex, multidimensional variations inherent in position embeddings.

\textbf{Search space}. We construct an extensive search space that captures the two identified non-uniformities, as detailed in Table~\ref{tbl:searchspace}. Our approach enables the selection of a distinct rescale factor for each individual dimension within the RoPE embedding. For simplicity in defining the search space, we optimize over $\lambda_i$ and $\hat{n}$, rather than directly exploring $\mathbb{I}(\hat{\lambda}_{i},\hat{n})$, where $\hat{\lambda}_{i}=1/\lambda_i$. As indicated in Table~\ref{tbl:searchspace}, the permissible range for $\lambda_i$ starts at 1.0 (corresponding to direct extrapolation) and extends up to $s\times1.25$ (which permits a broader interpolation range compared to PI), using increments of 0.01. Here, $s$ denotes the scaling factor for the intended extension of the context window.

The parameter $\hat{n}$ specifies how many of the earliest token positions preserve the standard RoPE embedding without any positional interpolation applied. In practice, we enable $\hat{n}$ to be selected from the set \{0, 1, 2, 4, 8, 12, 16, 20, 24, 28, 32, 64, 128, 256\}. Setting $\hat{n}=0$ results in every token position being subjected to the optimized rescaling factors.

\textbf{Evolution-based search}. The search space presented in Table~\ref{tbl:searchspace} encompasses a vast array of positional interpolation methods, making thorough exploration computationally demanding. For instance, when extending with $s=4\times$, the potential configurations number $400^{128/2}\times14=4\times10^{167}$. As the extension ratio increases, the complexity of the search space escalates exponentially. To manage this, we employ evolution search~\cite{guo2020single} and propose two additional optimization strategies to significantly enhance search efficiency. The comprehensive search process is outlined in Algorithm~\ref{alg:search}.

\textit{Optimized initial population generation}. Rather than creating an initial population of $P$ rescale factors purely at random, we explicitly include the three RoPE rescale factors associated with PI, NTK, and YaRN as members of the starting population. The other $P$-3 individuals are produced by introducing random mutations to these three rescale factors with a probability $p$.

\textit{Monotonically non-decreasing constraint}. Once the initial population has been produced, we evaluate each candidate by measuring its LLM perplexity. This involves applying the designated RoPE rescale factors within the target LLM, then assessing the perplexity on the input $\mathbf{X}$. The highest-scoring $k$ individuals are selected to serve as parents for the next generation. Nonetheless, because the search space is so expansive, basic mutation and crossover strategies can result in frequent evaluation of suboptimal candidates, incurring excessive perplexity computations. This inefficiency becomes especially pronounced for large $L'$, as computing each perplexity value is computationally intensive.

To tackle this issue, we enforce a non-decreasing monotonicity condition on the set of RoPE rescaled factors that are sampled: $\lambda_i \le \lambda_{i+1}$. Only those RoPE parameterizations that fulfill this condition are utilized for evaluating perplexity on the LLM, which markedly curtails the computational overhead of the search process. In detail, this means that $\lambda_i$ must monotonically increase with respect to the RoPE dimension index ($i$ = 0, ..., 63). The rationale for this dimensional monotonicity stems from NTK theory~\cite{jacot2018neural,tancik2020fourier,ntk}, which argues that in the lower-frequency, higher-indexed dimensions more interpolation (i.e., larger $\lambda_i$) is permissible, while in the higher-frequency, lower-indexed dimensions less interpolation is preferable (i.e., smaller $\lambda_i$).

\begin{algorithm}[t]
	\small
	\caption{The search algorithm}
	\textbf{Input:} target LLM, input samples $\mathbf{X}$,   population size $P$, mutation size $N_1$, crossover size $N_2$, max iterations $\mathcal{T}$, mutate probability $p$\\

	\begin{algorithmic}[1]
		\label{alg:search}
		\STATE Top-k $\gets \phi$;	
		\STATE $\text{P}_0 \gets \textit{Initialize\_population\_with\_optimization}(P, \mathbf{X}, p)$; 
		\FOR{$i = 1$ to $\mathcal{T}$}
		\STATE \textit{Compute\_perplexity} (LLM, $\text{P}_{i-1}$, $\mathbf{X}$);
		\STATE Top-k $\gets \textit{Update\_Topk}$ (Top-k, $\text{P}_{i-1}$);
		\STATE $\text{P}_{mutation} \gets \textit{Mutation\_with\_mono\_constraint}$ (Top-k, $N_1$, $p$); 
		\STATE $\text{P}_{crossover} \gets \textit{Crossover\_with\_mono\_constraint}$ (Top-k, $N_2$);
		\STATE $\text{P}_i \gets \text{P}_{mutation} \cup \text{P}_{crossover} \cup$ Top-k;
		\ENDFOR
		\STATE Output the entity in Top-k possessing the smallest perplexity value;
	\end{algorithmic}
\end{algorithm}

\textbf{8$\times$ extension without fine-tuning}. Through our evolutionary search approach, we successfully determine non-uniform RoPE rescale factors that retain important dimensions and positions, thereby reducing information loss caused by interpolation. As shown in Fig.\ref{fig:non-ft}, our method enables the expansion of LLaMA2’s context window from 4k to 32k tokens without any need for fine-tuning. By comparison, prior techniques like PI and non-uniform NTK and YaRN lead to a sharp increase in perplexity once the context window surpasses a 2$\times$ extension.

\subsection{Extending LLM Context Window to 2048K}
\label{sec:extendto2048k}

\textbf{Progressive extension to 2048k}. In this section, we present our approach for expanding the context window of pre-trained LLMs from the standard 4k up to more than 2048k tokens. Our results indicate that the application of non-uniform positional interpolation enables up to an 8$\times$ extension without the need for any additional fine-tuning. When attempting even more substantial extensions (such as 512$\times$), fine-tuning becomes essential. A conventional strategy involves identifying suitable RoPE rescaling parameters tailored to the 2048k context length before proceeding with fine-tuning. Nevertheless, this approach is hampered by the extremely high computational costs involved. Additionally, our observations suggest that achieving effective fine-tuning at such large extension scales poses significant difficulties (refer to Appendix for further details).

Fortunately, LongRoPE demonstrates strong performance with both base models and those that have been fine-tuned with extended context. Thus, we propose a resource-efficient and incremental strategy that reaches the desired 2048k target in only 1,000 fine-tuning steps, using training sequences of up to 256k tokens.

$\Diamond$ \textit{LongRoPE search for LLM extension to 256k context length}. Using LLaMA2 as a case study, we perform a search aimed at reaching context window sizes of 128k and 256k. At this stage, the enlargement factors are 32$\times$ and 64$\times$, respectively.

$\Diamond$ \textit{Fine-tuning to 256k}. In this phase, we adapt the pre-trained LLM for a 256k context window. The procedure involves initially fine-tuning LLaMA2 for 400 steps with RoPE rescaling set to 128k. Following this, we update the RoPE rescaling factors to target 256k on the obtained checkpoint and carry out a further 600 steps of fine-tuning. This approach has been found to be more efficient compared to directly fine-tuning at the 256k context length.

$\Diamond$ \textit{Expanding fine-tuned extended LLM to 2048k utilizing LongRoPE search}. In the concluding step, an additional search is conducted on the 256k context length LLM that was previously fine-tuned. Through this process, we achieve an exceptionally large context window of 2048k, all without requiring any additional fine-tuning steps. This corresponds to a final extension ratio of $512\times$.

\textbf{Recovery of performance in shorter context windows}.  When increasing the context window to an extensive size of 2048k, there is a discernible decrease in performance for the original, shorter context range. This limitation is well-documented in the literature on positional interpolation~\cite{pi}, which compresses the position embeddings of the original context window into a disproportionately compact subspace, thereby diminishing the model’s efficacy. Specifically, utilizing a 512$\times$ expansion means the positional representations within the initial 4k window are highly compressed and densely packed.

To address this issue, we conduct an additional evolution search on the extended LLM, targeting the optimization of RoPE rescale factors specifically for shorter context lengths such as 4k and 8k. The upper limit for the searched $\lambda$ is decreased because shorter contexts necessitate less positional interpolation. At inference time, the LLM adaptively updates the relevant RoPE rescale factors accordingly.

 \section{Experiments}

\label{sec:exp}
\subsection{Setup}
\textbf{Evaluation Tasks and Models}. We integrate LongRoPE into both LLaMA2-7B and Mistral-7B models, assessing their efficacy across three main dimensions: (1) the perplexity observed when processing lengthy documents with extended-context LLMs; (2) performance on the Passkey retrieval task, which assesses the model’s capability to identify a straightforward passkey within a large volume of distractor content; and (3) established LLM benchmarks evaluated under a standard 4096-token context window.

\textbf{Fine-tuning}. For LLaMA2, fine-tuning is performed using a learning rate of 2e-5 that decays linearly, with a total global batch size of 32. The model undergoes 400 training steps on the Redpajama~\cite{redpajama} dataset, which is partitioned into 128k-length chunks, each marked with BOS and EOS tokens at the boundaries. Subsequent to this initial fine-tuning, an additional 600 training steps are conducted—using the resulting checkpoint—to extend the model's context window to 256k. The initial 128k context phase is run on 8 A100 GPUs leveraging the distributed training framework~\cite{cube}, while training at 256k context involves 16 A100 GPUs. For the Mistral model, we opt for a constant learning rate of 1e-6 and a global batch size of 64. Both the 128k and 256k variants are tuned following the YaRN~\cite{yarn} configuration, which prescribes 400 training steps over Together Computer's Long-Data Collections~\cite{mistraldata} with a sequence length of 16k. All Mistral training is performed on 4 A100 GPUs.

\textbf{Search}. When the target window size does not exceed 256k, our configuration is as follows: $P$=64, $N_1$=$N_2$=16, $p$=0.3, $\mathcal{T}$=40. In every iteration, we pick the best 32 candidates for mutation and crossover. Perplexity metrics are obtained from 5 randomly selected samples drawn from the PG19 validation set, each of which meets or exceeds the target context length. For target windows exceeding 512k, the sizes of the population, mutation, and crossover pools are reduced by half. In this larger window setting, perplexity evaluation is conducted using 3 randomly chosen examples from the Pile-Books3~\cite{pile} validation dataset.

\textbf{Baselines}. To extend the context window to 2048k, we applied fine-tuning to models configured with 128k and 256k tokens, respectively. The result is LongRoPE-2048k (ft=128k) for LLaMA2 and LongRoPE-2048k (ft=256k) for Mistral. Our evaluation contrasts these four variants against leading context window extension baselines, focusing on publicly available LLMs that have undergone post-positional interpolation fine-tuning through methods such as PI, NTK, and YaRN. The benchmark models considered include Together-32k~\cite{together}, Code LLaMA~\cite{codellama}, LongLoRA-full-FT-100k~\cite{longlora}, YaRN-LLaMA, and YaRN-Mistral~\cite{yarn}.

\subsection{Main Results}

\label{sec:mainresult}
\textbf{Long sequence language modeling up to 256k tokens}. Our initial evaluation involves benchmarking our method against leading extended LLMs over sequence lengths reaching 256k tokens. To highlight our approach's generalizability, we utilize test splits from two datasets: Proof-pile~\cite{pg19} and PG19~\cite{pile}. Perplexity is measured across multiple context lengths, applying a sliding window size of 256. For the PG19 corpus, the evaluation incorporates all 100 test set documents. For Proof-pile, we adhere to the protocol outlined by YaRN~\cite{yarn}, randomly selecting 10 samples, each containing at least 128k tokens.

Table~\ref{tbl:proof-pile} and Table~\ref{tbl:pg19} present a comparison of the perplexity scores for LLaMA2 and Mistral models, each expanded using various interpolation strategies, evaluated on Proof-pile and PG19, respectively. Two main findings emerge: \textbf{(1)} the perplexity of our extended models consistently declines as the evaluation length increases from 4k to 256k, indicating that they are able to exploit longer contexts effectively. \textbf{(2)} Remarkably, even when the models operate over a context window that is 16$\times$ larger—a scenario generally associated with significant performance degradation at shorter sequence lengths—our LongRoPE-2048k models surpass state-of-the-art counterparts within the 256k context length regime.

\textbf{Modeling extremely long sequences beyond 2000k}. To assess performance on very large text samples, we utilize the Books3~\cite{pile} dataset. For efficient evaluation, we choose a random subset of 20 books, each with a length greater than 2048k, and apply a 256k sliding window during testing.

As illustrated in Table~\ref{tbl:books3}, LongRoPE enables both LLaMA2-7B and Mistral-7B models to reach context lengths of up to 2048k, while delivering perplexity scores that match or surpass those of baseline models at shorter sequence lengths ranging from 8k to 128k. Distinct performance patterns emerge between LLaMA2 and Mistral when extended to 2048k: Mistral demonstrates superior results to baseline models at shorter contexts, but its perplexity increases to over 7 when the window exceeds 256k tokens. LLaMA2 behaves as anticipated, displaying reduced perplexity with larger context windows, though slight increases are observable at the 1024k and 2048k marks. Furthermore, LongRoPE-2048k applied to LLaMA2 yields enhanced results when fine-tuned with a 256k window size versus 128k, attributable to a less aggressive secondary extension ratio (specifically 8$\times$ compared to 16$\times$). On the other hand, Mistral attains better outcomes with a fine-tuning context of 128k. This discrepancy is mainly due to our adoption of YaRN's approach for Mistral, which uses a 16k training length during both 128k and 256k fine-tuning procedures—constraining Mistral's capacity for subsequent context window extension after fine-tuning.

\textbf{Passkey retrieval}. Next, we investigate the practical context window length for generative tasks. To this end, we employ a synthetic passkey retrieval benchmark as outlined by ~\cite{passkey}. In this evaluation, the model's objective is to identify a random passkey—a five-digit numeric code—embedded within an extensive document. The specific prompt template utilized for this experiment is described in the appendix. For assessment, we conduct 10 rounds of the task, each time randomly positioning the passkey at a location drawn uniformly from the available context window.

As illustrated in Fig.~\ref{fig:passkey}, retrieval accuracy was compared against baseline models. While the accuracy of existing methods declines sharply to zero when the input length exceeds 128k tokens, our LongRoPE-LLaMA2-2048k (ft=256k) demonstrates consistently high retrieval accuracy ($\ge$90\%) across a wide range from 4k up to 2048k tokens, despite the inherently difficult task of passkey retrieval from sequences containing millions of tokens. For LongRoPE-Mistral-2048k (ft=128k), accuracy remains perfect (100\%) until 1800k tokens, after which it decreases to 60\% at 2048k. This drop is consistent with the pattern observed in Table~\ref{tbl:books3}, where perplexity exhibits a mild increase at the 2048k length.

\textbf{Established benchmarks under the original context window}. We assess the performance of LongRoPE-2048k models within their initial context window by utilizing the Hugging Face Open LLM Leaderboard~\cite{huggingfaceboard} for both zero-shot and few-shot scenarios. Specifically, our evaluation comprises 25-shot ARC-Challenge~\cite{allenai:arc}, 10-shot HellaSwag~\cite{zellers2019hellaswag}, 5-shot MMLU~\cite{mmlu}, and 0-shot TruthfulQA~\cite{lin2021truthfulqa}.

Table~\ref{tbl:commonsense} illustrates that our models perform on par with the initial benchmark—which was intended for shorter context lengths—and, in the case of TruthfulQA, surpass the original Mistral by an increase of +0.5\%. LongRoPE-LLaMA2-2048k, despite being fine-tuned at 256k and exhibiting a modest decrease in performance, maintains results that are generally acceptable across the evaluated tasks.

\subsection{Ablation Results}

\textbf{Effectiveness of the second positional interpolation}. Within our progressive extension approach, we employ our search algorithm to perform a second stage of non-uniform positional interpolation on the fine-tuned extended LLMs. To assess the impact of this technique, we conduct experiments utilizing our fine-tuned LLaMA2-256k model. The model is further extended to 512k, 1024k, and 2048k using both PI and YaRN methods. Results presented in Table~\ref{tbl:secondsearch} demonstrate that our approach to non-uniform positional interpolation maintains stable perplexity levels. In comparison, perplexity values associated with PI and YaRN rise sharply as the extension ratio increases.

\textbf{Effectiveness of recovery at shorter context lengths.} To address degraded performance at shorter contexts, we modify the RoPE scaling factors used in LongRoPE-2048k by applying our search methodology. In this process, we constrain the upper limit of allowable scaling during the search to promote reduced interpolation at short context lengths of 4k and 8k. Table~\ref{tbl:shortperformance} presents perplexity outcomes for LongRoPE-LLaMA2-2048k on the Proof-pile dataset for both 4k and 8k contexts, as well as average accuracy over a range of LLM benchmarks. The findings indicate that these modifications yield substantial gains in short-context performance.

\textbf{Analysis on the two forms of non-uniformities}. To further understand the impact of each type of non-uniformity, we conduct an ablation study examining their respective contributions to overall model performance. Our investigation involves two primary experimental setups: (i) We expand LLaMA2-7B to sequence lengths of 16k and 32k using three strategies—PI, modifying only the RoPE dimension, and modifying both non-uniform aspects; (ii) We take our already fine-tuned LLaMA2 at 256k length and extend it to 2048k, employing the same range of methods. Perplexity measurements are obtained without additional fine-tuning. According to Table~\ref{tbl:searchdimension}, introducing non-uniformity into the RoPE dimension leads to notable decreases in perplexity, outperforming PI’s standard linear interpolation. Incorporating non-uniformity in token position also enhances results for shorter contexts (16k and 32k), yet this advantage diminishes at the much larger 2048k scale—likely attributable to the excessive sequence length, where retaining only the starting tokens without interpolation yields little benefit. Investigation of these dynamics at extreme lengths will be addressed in future studies.

\section{Related Works}

Besides position interpolation-based techniques, this section reviews alternative related methods.

\textbf{Retrieval-based} techniques employ an external memory component to store extensive past contexts, leveraging retrieval mechanisms to access pertinent documents during inference~\cite{longllama,wang2023augmenting,borgeaud2022improving}. Such approaches generally require substantial architectural changes to LLMs. By comparison, our method is more lightweight, requiring only minimal changes to positional embeddings. Furthermore, our approach extends beyond retrieval tasks and effectively addresses a broader array of long-context problems, including tasks such as long document summarization and few-shot learning.

\textbf{Attention-based context window extensions}. In addition to positional embedding interpolation, certain studies have explored the extension of input context by retaining the standard LLM context window size and altering the attention mechanism~\cite{han2023lminfinite, streamingllm,parallelwindow}. The primary strategy involves introducing innovative attention masks to address the problem of attention explosion when handling new positions. Such approaches are complementary to methods based on positional interpolation.

\textbf{Fine-tuning based approaches} explore strategies for adapting pre-trained LLMs with altered position embeddings to accommodate longer input sequences. Methods such as Code LLaMA~\cite{codellama}, LLaMA2 Long~\cite{xiong2023effective}, and ScaledRoPE~\cite{liu2023scaling} implement this by adopting a substantially larger base value for RoPE, subsequently fine-tuning on the intended sequence length. In contrast, our approach allows for flexible adaptation to multiple target lengths and scales to over 2M tokens. With the increasing resource demands of fine-tuning at extensive context lengths (e.g., above 128k tokens), solutions like LongLoRA~\cite{longlora} and PoSE~\cite{zhu2023pose} have been introduced to alleviate GPU usage. Our approach is complementary to these efficient fine-tuning strategies.

\section{Conclusion}

In this study, we introduce LongRoPE, a technique that significantly increases the context window of LLMs to an extraordinary 2048k tokens, while preserving performance within their initial, shorter context boundaries. Our approach leverages two distinct non-uniformities in the RoPE positional encoding through a streamlined evolutionary search. This methodology yields dual advantages: it supplies effective initialization for subsequent fine-tuning, and allows for an 8$\times$ context window extension even in the absence of fine-tuning. Furthermore, we develop a stepwise expansion strategy whereby LLMs fine-tuned for 256k tokens are utilized to achieve a 2048k context range without the need for additional fine-tuning. Comprehensive experimental results confirm the efficacy of LongRoPE. We anticipate that our LongRoPE-2048k models will unlock a range of new applications requiring extended contexts and will stimulate further research in this area.

\section*{Broader Impacts} The research detailed in this paper aims to contribute to progress in the domain of Machine Learning. While the implications of this work may extend to various societal contexts, we do not identify any particular consequences that require explicit mention at this time.

	\balance

\end{document}